\chapter{Proposta}
\label{cap:proposta}

\section{Contextualização e Motivação}
\label{sec:P-contexto}

O aprendizado profundo consolidou-se como abordagem dominante em tarefas de
classificação e segmentação de imagens médicas. Esse sucesso, contudo, apoia-se
em uma premissa que a prática clínica raramente satisfaz: a disponibilidade de
grandes conjuntos de dados anotados. Em imagens médicas, a anotação exige um
especialista, é onerosa e lenta, e frequentemente esbarra em restrições éticas
e de privacidade que limitam o compartilhamento entre instituições.

Essa escassez tem uma consequência metodológica pouco discutida. O procedimento
padrão para decidir quando interromper o treinamento --- o \emph{early
stopping} --- monitora o erro em um conjunto de validação separado e interrompe
o ajuste quando esse erro deixa de diminuir \cite{prechelt1998}. O
procedimento é sólido, mas transfere para o conjunto de validação toda a
responsabilidade pelo diagnóstico: se ele é pequeno, desbalanceado ou pouco
representativo, a estimativa do momento de parada torna-se ruidosa, e o
desenvolvedor não dispõe de nenhum sinal alternativo para confrontá-la. Em
outras palavras, justamente no regime de poucos dados --- aquele em que o
\emph{overfitting} é mais provável --- o instrumento de medida é o menos
confiável.

Há, entretanto, uma fonte de informação sobre o estado do treinamento que
permanece pouco explorada e que não depende de rótulo algum: os próprios pesos
da rede. A memorização de exemplos deixa marcas na distribuição de valores dos
parâmetros, e trabalhos consolidados mostram que redes profundas são capazes de
ajustar rótulos completamente aleatórios \cite{zhang2017}, memorizando-os por um
regime de aprendizado distinto daquele em que extraem padrões
generalizáveis \cite{arpit2017}. Se esses dois regimes deixam assinaturas
distinguíveis na estrutura dos pesos, então medidas de complexidade estatística
aplicadas a esses pesos podem, em princípio, sinalizar a transição entre eles
sem consultar nenhum dado rotulado.

\section{Problema de Pesquisa}
\label{sec:P-problema}

É possível detectar o início do \emph{overfitting} de uma rede neural profunda
a partir exclusivamente da complexidade estatística de seus pesos, sem recorrer
a um conjunto de validação rotulado?

\section{Hipótese}
\label{sec:P-hipotese}

A hipótese central deste trabalho é a de que a transição do regime de
aprendizado generalizante para o regime de memorização produz uma mudança
mensurável na irregularidade estrutural dos pesos da camada densa, e que essa
mudança \emph{antecede} a degradação observável no conjunto de validação.

Operacionalmente, espera-se que, em execuções nas quais o \emph{overfitting} é
induzido de forma controlada, ao menos um dos indicadores considerados apresente
inflexão em uma época anterior àquela em que a acurácia de validação começa a
cair de modo sustentado, e que essa inflexão não ocorra --- ou ocorra de forma
substancialmente atenuada --- em execuções de controle sem indução. A hipótese
é falsificável: caso a inflexão ocorra simultaneamente ou posteriormente à
degradação da validação, ou caso ocorra com a mesma intensidade nas execuções
de controle, o indicador não tem valor como aviso antecipado e isso deve ser
reportado como tal.

\section{Objetivos}
\label{sec:P-objetivos}

\subsection{Objetivo Geral}

Avaliar a complexidade estatística dos pesos de redes neurais profundas como
indicador de \emph{overfitting} independente do conjunto de validação,
entregando ao desenvolvedor um instrumento de diagnóstico acoplável ao laço de
treinamento.

\subsection{Objetivos Específicos}

\begin{enumerate}
    \item Implementar, em um módulo único e validado numericamente, a
          complexidade LMC, a entropia amostral unidimensional (SampEn) e a
          entropia amostral bidimensional (SampEn2D), estabelecendo-o como
          fonte única de cálculo para todo o trabalho.
    \item Formalizar a leitura dos pesos da camada densa como objeto de
          medida, fixando a ordem de achatamento da matriz de pesos e
          justificando-a experimentalmente.
    \item Definir operacionalmente o início do \emph{overfitting} e construir
          um protocolo de indução controlada por corrupção de rótulos, com
          execuções de controle pareadas.
    \item Implementar um monitor que calcule os indicadores a cada época
          durante o treinamento e emita um estado do aprendizado em tempo real.
    \item Quantificar a antecedência de cada indicador em relação à degradação
          da validação, com múltiplas sementes e teste de significância.
    \item Caracterizar os limites do indicador quanto ao nível de corrupção e
          quanto à arquitetura da rede, reportando explicitamente os regimes em
          que ele falha.
\end{enumerate}

\section{Justificativa: o ponto de vista do desenvolvedor}
\label{sec:P-justificativa}

A justificativa deste trabalho é deliberadamente formulada do ponto de vista de
quem treina a rede, e não apenas do ponto de vista teórico.

\subsection{A dependência atual dos dados rotulados}

Para saber se uma rede aprendeu ou apenas memorizou, o desenvolvedor depende
hoje inteiramente dos dados. É preciso construir conjuntos de treino, validação
e teste que sejam bem particionados, corretamente rotulados e representativos da
população de interesse. Cada um desses requisitos custa tempo e dinheiro, e em
imagens médicas frequentemente esbarra na indisponibilidade de especialistas
para anotação. O diagnóstico do treinamento é, portanto, uma função do orçamento
de anotação --- não uma propriedade que se possa ler diretamente do modelo.

\subsection{O que muda com um indicador calculado sobre os pesos}

Os indicadores investigados neste trabalho são calculados exclusivamente a
partir dos pesos da rede. Não consomem imagem, não consomem rótulo e não
dependem de nenhuma partição. Se forem capazes de sinalizar a transição para o
regime de memorização, o desenvolvedor passa a dispor de uma segunda leitura do
treinamento, obtida por um caminho independente daquele do conjunto de
validação. As duas leituras podem então ser confrontadas: quando concordam, a
decisão de parada ganha confiança; quando discordam, há indício de que a
validação não é representativa.

\subsection{Casos de uso concretos}

\begin{enumerate}
    \item \textbf{Decidir quando parar.} Um aviso emitido antes da degradação
          observável na validação amplia a janela de decisão do desenvolvedor,
          em vez de constatar a perda depois de consumada.
    \item \textbf{Economizar computação.} Detectar que a rede deixou de
          aprender evita épocas de treinamento que não trarão ganho, com
          impacto direto em custo de GPU.
    \item \textbf{Comparar arquiteturas sem ampliar a anotação.} Como a medida
          não consome dados rotulados, a comparação entre arquiteturas
          candidatas não exige a ampliação do conjunto de validação.
    \item \textbf{Diagnosticar com validação pequena.} Em cenários nos quais o
          conjunto de validação é pequeno demais para ser confiável, um
          indicador independente oferece evidência adicional --- exatamente o
          cenário típico em imagens médicas.
\end{enumerate}

\subsection{Entrega prática}

A contribuição não se esgota na constatação empírica. Entrega-se um monitor
implementado em PyTorch, acoplável ao laço de treinamento de terceiros, que
calcula os indicadores a cada época e devolve um estado do aprendizado. O custo
computacional é baixo, pois a medida incide apenas sobre a camada densa.

\section{Visão Geral da Metodologia}
\label{sec:P-metodologia}

O trabalho experimental organiza-se em torno da comparação pareada entre
execuções com e sem \emph{overfitting} induzido. A indução é feita por corrupção
controlada de uma fração dos rótulos do conjunto de treino, procedimento que
força a rede a memorizar exemplos inconsistentes \cite{zhang2017}. Cada execução
com corrupção possui uma execução de controle correspondente, idêntica em
arquitetura, partição e semente de inicialização, diferindo apenas pela
presença da corrupção --- de modo que a semente que governa o ruído é separada
daquela que governa a partição e a inicialização.

Durante o treinamento, um monitor calcula a cada época a LMC, a SampEn e a
SampEn2D sobre os pesos da camada densa. Ao final, compara-se a época em que
cada indicador inflete com a época em que a acurácia de validação passa a cair
de forma sustentada, definida como âncora. A diferença entre as duas é a
antecedência, positiva quando o indicador avisa antes. A significância é
avaliada por teste não paramétrico entre os grupos com e sem corrupção,
escolhido pelo tamanho reduzido das amostras.

O detalhamento reprodutível --- repositório de imagens, arquiteturas,
hiperparâmetros, partição congelada e critério de inflexão --- constitui o
capítulo de metodologia do documento final.

\section{Contribuição Esperada}
\label{sec:P-contribuicao}

Espera-se contribuir em três frentes. Do ponto de vista científico, caracterizar
empiricamente se e quando medidas de complexidade estatística sobre pesos
antecipam o \emph{overfitting}, incluindo a delimitação honesta dos regimes em
que não o fazem. Do ponto de vista metodológico, formalizar a leitura dos pesos
de uma camada densa como objeto de medida de complexidade, questão até aqui
tratada de modo implícito na literatura. Do ponto de vista prático, disponibilizar
um monitor documentado e testado, utilizável por terceiros independentemente
deste trabalho.
